\section*{अध्याय \ref{ch:g7:earthquakes} --- भूकंप}

\begin{solution}{exo:g7:earthquakes:1}
(क) \omterm{def:g7:earthquakes:focus}{भूकंप-मूल}; (ख) \omterm{def:g7:earthquakes:intensity}{तीव्रता}; (ग) \omterm{def:g7:earthquakes:fault}{भ्रंश}; (घ) \omterm{def:g7:earthquakes:focus}{अधिकेंद्र}; (ङ) \omterm{def:g7:earthquakes:magnitude}{परिमाण}; (च)
\omterm{def:g7:earthquakes:seismometer}{भूकंपमापी}।
\end{solution}

\begin{solution}{exo:g7:earthquakes:2}
(क) गाँव पर ही: \omterm{def:g7:earthquakes:focus}{अधिकेंद्र} सतह का वह बिंदु है जो \omterm{def:g7:earthquakes:focus}{भूकंप-मूल} के ठीक ऊपर है।
(ख) \qty{30}{km}, सीधे नीचे। (ग) आमतौर पर कम: दूरी के साथ काँपना कमज़ोर पड़ता है।
(यह ज़मीन पर भी निर्भर करता है: नरम \omterm{def:g3:stones-pebbles-sand:sand}{कीचड़} पर दूसरा गाँव फिर भी ज़ोर से काँप
सकता है।)
\end{solution}

\begin{solution}{exo:g7:earthquakes:3}
(क) \omterm{def:g7:earthquakes:magnitude}{परिमाण} (एक दशमलव स्थान वाली संख्या)। (ख) \omterm{def:g7:earthquakes:intensity}{तीव्रता} (एक जगह पर जो दिखा)।
(ग) \omterm{def:g7:earthquakes:intensity}{तीव्रता} (रोमन अंक, दो जगहों पर अलग)। (घ) \omterm{def:g7:earthquakes:magnitude}{परिमाण}: छोड़ी गई ऊर्जा, \omterm{def:g7:earthquakes:magnitude}{परिमाण} का एक
क़दम ज़्यादा।
\end{solution}

\begin{solution}{exo:g7:earthquakes:4}
(क) ग़लत: वैज्ञानिक \omterm{def:g4:volcanoes-earthquakes-first-look:earthquake}{भूकंपों} की भविष्यवाणी नहीं कर सकते; वे केवल यह जानते हैं कि
कौन-से \omterm{def:g7:earthquakes:fault}{भ्रंश} किसी दिन फिर टूटेंगे। (ख) ग़लत: खुले \omterm{def:g2:water-on-earth:ocean}{महासागर} में \omterm{def:g7:earthquakes:tsunami}{सुनामी} शायद ही कभी
एक मीटर से ज़्यादा ऊँची होती है; वह \omterm{def:g1:landscapes-around-us:water}{तट} के पास ऊँची होती है। (ग) सही। (घ) सही:
\omterm{def:g7:earthquakes:magnitude}{परिमाण}~6 या उससे ज़्यादा के 92~\% \omterm{def:g4:volcanoes-earthquakes-first-look:earthquake}{भूकंप} किसी \omterm{def:g7:plate-tectonics:plate-boundary}{प्लेट सीमा} से
\qty{300}{km} से कम दूरी पर आते हैं।
\end{solution}

\begin{solution}{exo:g7:earthquakes:5}
(क) \omterm{def:g7:earthquakes:magnitude}{परिमाण} 6 से 6.9 के लगभग 127 प्रति वर्ष; \omterm{def:g7:earthquakes:magnitude}{परिमाण} 7 से 7.9 के लगभग 13।
(ख) $126.7 \div 13.12 \approx 9.7$: लगभग 10~गुना कम। (ग)
$13.12 + 0.76 + 0.04 = 13.92$: लगभग \textbf{14} प्रति वर्ष। (घ)
$0.04 \times 50 = 2$: 50~वर्षों में दो (सुमात्रा 2004 और जापान 2011)।
\end{solution}

\begin{solution}{exo:g7:earthquakes:6}
(क) एक क़दम: 32~गुना। (ख) दो क़दम: $32 \times 32 = 1024$, लगभग
1000~गुना। (ग) हर क़दम तरंगों को 10~गुना बड़ा करता है: $10 \times 10 =
100$~गुना बड़ी।
\end{solution}

\begin{solution}{exo:g7:earthquakes:7}
(क) $8 \times 5 = \qty{40}{km}$। (ख) $8 \times 25 = \qty{200}{km}$। (ग)
$120 \div 8 = \qty{15}{s}$। (घ) एक स्टेशन केवल दूरी देता है, दिशा नहीं: उस दूरी
वाले सारे बिंदु स्टेशन के चारों ओर एक वृत्त बनाते हैं, और \omterm{def:g7:earthquakes:focus}{अधिकेंद्र} उस पर कहीं भी
हो सकता है।
\end{solution}

\begin{solution}{exo:g7:earthquakes:8}
(क) V। (ख) II। (ग) VIII। (घ) \omterm{def:g7:earthquakes:intensity}{तीव्रता} जगह-जगह बदलती है: रिपोर्ट (ग) \omterm{def:g7:earthquakes:focus}{अधिकेंद्र}
के पास की, या नरम ज़मीन पर बसी किसी जगह से आई है; रिपोर्ट (क) ज़्यादा दूर की, या
ठोस \omterm{def:g3:stones-pebbles-sand:rock}{चट्टान} पर बसी किसी जगह से। \omterm{def:g4:volcanoes-earthquakes-first-look:earthquake}{भूकंप} का फिर भी केवल एक ही \omterm{def:g7:earthquakes:magnitude}{परिमाण} है।
\end{solution}

\begin{solution}{exo:g7:earthquakes:9}
(क) $7031 - 5561 = 1470$। (ख) $1470 \div 7031 \approx 0.209$: लगभग
\textbf{21~\%}। (ग) धँसती \omterm{def:g7:plate-tectonics:plate}{प्लेट} \omterm{def:g6:inside-the-earth:layers}{प्रावार} में गहरे तक नीचे जाती है। वह अब भी इतनी
कठोर है कि टूट सके, इसलिए \omterm{def:g4:volcanoes-earthquakes-first-look:earthquake}{भूकंप} उसके अंदर, सैकड़ों किलोमीटर नीचे, शुरू हो
सकते हैं।
\end{solution}

\begin{solution}{exo:g7:earthquakes:10}
(क) $2500 \div 800 = 3.125$~घंटे: 3~घंटे और $0.125 \times 60 = 7.5$
मिनट, लगभग 3~घंटे 8~मिनट। (ख) 3~घंटे~8~मिनट $-$ 2~घंटे~17~मिनट $=$
51~मिनट बाद। (ग) लहर को पहुँचने में घंटों लगे, जबकि \omterm{def:g2:water-on-earth:ocean}{समुद्र} के नीचे के बड़े \omterm{def:g4:volcanoes-earthquakes-first-look:earthquake}{भूकंप} का
पता कुछ ही मिनटों में चल जाता है: चेतावनी लहर आने से बहुत पहले लोगों को ऊँची ज़मीन
पर भेज सकती थी।
\end{solution}

\begin{solution}{exo:g7:earthquakes:11}
(क) $400 \div 800 = 0.5$~घंटा: 30~मिनट। (ख) $20 \div 40 = 0.5$~घंटा:
30~मिनट। (ग) 60~मिनट। (घ) आख़िरी \qty{20}{km} में उतना ही समय लगता है जितना
पहले \qty{400}{km} में: उथले पानी में लहर बहुत ज़्यादा धीमी हो जाती है। (ङ)
लहर का अगला हिस्सा धीमा हो जाता है जबकि उसके पीछे का पानी अब भी तेज़ी से आता है:
पानी जमा होता जाता है, और लहर ऊँची हो जाती है।
\end{solution}

\begin{solution}{exo:g7:earthquakes:12}
(क) तीन क़दम: $32 \times 32 \times 32 = 32\,768$, लगभग
\textbf{33\,000}~गुना। (ख) $32\,768 \div 127 \approx 258$: दुनिया के परिमाण-6
वाले सभी \omterm{def:g4:volcanoes-earthquakes-first-look:earthquake}{भूकंपों} के लगभग 260~वर्ष। (ग) लगभग सारी ऊर्जा गिने-चुने सबसे बड़े \omterm{def:g4:volcanoes-earthquakes-first-look:earthquake}{भूकंप}
छोड़ते हैं: \omterm{def:g7:earthquakes:magnitude}{परिमाण}~9 का एक अकेला \omterm{def:g4:volcanoes-earthquakes-first-look:earthquake}{भूकंप} परिमाण-6 वाले \omterm{def:g4:volcanoes-earthquakes-first-look:earthquake}{भूकंपों} की सदियों पर भारी
पड़ता है।
\end{solution}

\begin{solution}{pb:g7:earthquakes:1}
\textbf{1.} A: $45 - 25 = 20$~सेकंड। B: $29 - 16 = 13$~सेकंड। C: $34 - 19 =
15$~सेकंड।

\textbf{2.} स्टेशन B: उसकी देरी सबसे कम है, और देरी जितनी कम, \omterm{def:g4:volcanoes-earthquakes-first-look:earthquake}{भूकंप} उतना पास।
(उसकी तेज़ तरंगें भी सबसे पहले पहुँचीं।)

\textbf{3.} A: $8 \times 20 = \qty{160}{km}$। B: $8 \times 13 =
\qty{104}{km}$। C: $8 \times 15 = \qty{120}{km}$।

\textbf{4.} A से \qty{160}{km} दूर के सभी बिंदु: A पर केंद्रित
\qty{160}{km} त्रिज्या का एक वृत्त। \omterm{def:g7:earthquakes:focus}{अधिकेंद्र} उस पर कहीं है।

\textbf{5.} A: $160 \div 20 = 8$~वर्ग। B: $104 \div 20 = 5.2$~वर्ग।
C: $120 \div 20 = 6$~वर्ग।

\textbf{6.} E, O से $3 + 6 = 9$~वर्ग पूर्व और 14.5~वर्ग उत्तर में है, तीनों
स्टेशनों के बीच \omterm{def:g2:water-on-earth:ocean}{समुद्र} में; A और B के वृत्त उससे गुज़रते हैं।

\textbf{7.} \omterm{def:g2:water-on-earth:ocean}{समुद्र} के नीचे। \omterm{def:g2:water-on-earth:ocean}{समुद्र} की तली की अचानक गति \omterm{def:g7:earthquakes:tsunami}{सुनामी} पैदा कर सकती है:
तटों और द्वीप को चेतावनी देनी होगी।

\textbf{8.} \omterm{def:g7:earthquakes:magnitude}{परिमाण} का एक क़दम: लगभग 32~गुना ज़्यादा ऊर्जा।

\textbf{9.} $12 \times 20 = \qty{240}{km}$।

\textbf{10.} $240 \div 800 = 0.3$~घंटा, और $0.3 \times 60 = 18$~मिनट।

\textbf{11.} \omterm{def:g4:volcanoes-earthquakes-first-look:earthquake}{भूकंप} के 5~मिनट बाद।

\textbf{12.} $18 - 5 = 13$: पोर्ट हेल के लोगों के पास ऊँची ज़मीन तक पहुँचने के लिए
लगभग \textbf{13~मिनट} हैं। (द्वीप के पास के उथले पानी में लहर धीमी हो जाती है,
जिससे थोड़ा और समय मिलता है, पर किसी को इस पर भरोसा नहीं करना चाहिए: वे तुरंत
निकल जाते हैं।)
\end{solution}
